\documentclass[11pt,a4paper]{article}
\usepackage{fontspec}
\usepackage[english,russian]{babel}
\setmainfont{Times New Roman}
\setsansfont{DejaVu Sans}
\setmonofont{DejaVu Sans Mono}[Scale=0.82]
\usepackage{amsmath,amssymb}
\usepackage[margin=2cm]{geometry}
\usepackage{xcolor}
\usepackage[newfloat]{minted}
\usepackage{booktabs}
\usepackage{tabularx}
\usepackage{enumitem}
\usepackage[ruled,linesnumbered]{algorithm2e}
\usepackage[colorlinks=true,linkcolor=blue!50!black,urlcolor=blue!50!black]{hyperref}
\setlist{nosep,leftmargin=*}

\definecolor{codebg}{RGB}{248,248,248}
\setminted{fontsize=\small, bgcolor=codebg, linenos, numbersep=6pt, breaklines, breakanywhere,
           frame=lines, framesep=4pt, style=friendly, tabsize=4}
\renewcommand{\theFancyVerbLine}{\textcolor{gray}{\tiny\arabic{FancyVerbLine}}}
\newcommand{\codefile}[1]{\inputminted{python}{#1}}
\newcommand{\code}[1]{\texttt{#1}}
\newcommand{\abar}{\bar\alpha}
\newcommand{\E}{\mathbb{E}}

\title{Реимплементация DimExDAM:\\ генерация 3D-микроструктур кернов по 2D-срезам\\[4pt]
\large Описание кода, формулы и полный листинг}
\author{Unarov Aital}
\date{Октябрь 2026}

\begin{document}
\maketitle
\tableofcontents
\newpage

%=====================================================================
\section{Обзор}

Проект реализует метод \emph{DimExDAM} (Dimensionality Expansion Diffusion Adversarial Model) из статьи
Aouf, Laloy, Rogiers, De Vleeschouwer, «Generating Heterogeneous 3D Geological Microstructures from 2D
Images via a Stable Diffusion-Adversarial Model» (arXiv 2609.20358).

\textbf{Задача.} По трём ортогональным 2D-срезам бинарного керна (поры/твёрдая фаза) научиться
генерировать новые 3D-объёмы, статистически эквивалентные исходному материалу. Модель безусловная:
на каждый керн (вариант) обучается своя модель.

\textbf{Идея метода.} Гибрид диффузии и GAN (DDGAN): 3D-генератор за короткую цепочку из $T=8$ шагов
превращает шум в объём, а единственный 2D-дискриминатор на каждом шаге сравнивает пары срезов
сгенерированного объёма с зашумлёнными реальными 2D-срезами. Вместо MSE-лосса денойзинга используется
адверсариальный лосс, поэтому 3D ground truth не нужен.

\textbf{Важно.} Это собственная реимплементация по статье. Цикл обучения и часть гиперпараметров сверены с
кодом авторов (\url{https://github.com/40uf411/DimExDAM}), но архитектура генератора упрощена
(3D U-Net вместо NCSN++), а дискриминатор написан по описанию в статье. Отличия перечислены в
разделе~\ref{sec:diff}.

\subsection{Структура проекта}
\begin{center}
\begin{tabularx}{\textwidth}{@{}lX@{}}
\toprule
Файл & Назначение \\
\midrule
\code{src/dimexdam/data.py} & чтение \code{.mhd}, обучающие срезы, кропы, размытие, стратифицированная нарезка 3D$\to$2D \\
\code{src/dimexdam/diffusion.py} & расписание шума, прямой процесс, апостериорное распределение, сэмплинг \\
\code{src/dimexdam/models.py} & 3D-генератор (U-Net с AdaGN) и 2D-дискриминатор (spectral norm) \\
\code{src/dimexdam/train.py} & конфигурация, цикл обучения, сохранение/загрузка чекпойнтов \\
\code{src/dimexdam/metrics.py} & доля фазы, двухточечная функция $S_2(r)$, случайные подобъёмы \\
\code{scripts/generate\_samples.py} & обучение (если нет чекпойнта) и генерация $N$ объёмов, замеры времени и памяти \\
\code{scripts/calibrate\_threshold.py} & порог бинаризации по пористости 2D-срезов, серые объёмы \\
\code{scripts/view\_volume.py} & 3D-рендер, листалка срезов, экспорт в \code{.mhd} \\
\code{notebooks/dimexdam\_inference.ipynb} & ноутбук: обучение $\to$ инференс $\to$ сравнение с керном \\
\bottomrule
\end{tabularx}
\end{center}

\subsection{Поток данных}
\begin{enumerate}
  \item Керн $300^3$ (\code{data/Variant\_XX/model.mhd}) $\to$ три центральных среза $300\times300$
        (\code{ortho\_slices}). Остальной объём модель не видит.
  \item Обучение: случайные кропы $64\times64$ из срезов против стратифицированных срезов сгенерированных
        объёмов $64^3$, на каждом шаге диффузии.
  \item Инференс: шум $\to$ 8 проходов генератора $\to$ «серый» объём в $[-1,1]$ $\to$ бинаризация
        (порог 0 или калиброванный).
  \item Результат: \code{outputs/Variant\_XX/samples.npy}, массив $(N,64,64,64)$, \code{uint8},
        $0$ = поры, $1$ = твёрдая фаза.
\end{enumerate}

%=====================================================================
\section{Данные (\code{data.py})}

\subsection{Формат и нормировка}
Керн хранится в формате MetaImage: заголовок \code{.mhd} и сжатые данные \code{.raw.gz}. Размеры в
заголовке указаны в порядке $X\,Y\,Z$ (быстрее всего меняется $X$), поэтому массив читается как
$(Z,Y,X)$. Значения бинарные, $v\in\{0,1\}$, где $0$ — поры. Для сети они переводятся в $\{-1,1\}$:
\[
  x = 2v - 1,\qquad \text{поры}\mapsto -1,\quad \text{твёрдая фаза}\mapsto +1 .
\]

\subsection{Обучающие изображения}
Срез $p\in\{0,1,2\}$ перпендикулярен оси $p$ и берётся в центре керна:
\[
  I_0 = V[150,:,:],\qquad I_1 = V[:,150,:],\qquad I_2 = V[:,:,150].
\]
Это весь обучающий набор: по одному изображению $300\times300$ на плоскость, как в статье.

\subsection{Кропы и low-pass curriculum}
\code{SliceSampler} выдаёт $n$ кропов $L\times L$ ($L=64$): для каждого случайно выбирается плоскость $p$
и положение кропа. Часть кропов размывается гауссовым ядром (curriculum, раздел~\ref{sec:curr}).
Ядро сепарабельное, с отражающим паддингом:
\[
  k_i = \frac{e^{-i^2/(2\sigma^2)}}{\sum_{j=-r}^{r} e^{-j^2/(2\sigma^2)}},\qquad i=-r,\dots,r,\quad r=\max(1,\lfloor 3\sigma\rfloor).
\]

\subsection{Стратифицированная нарезка 3D $\to$ 2D}
Из каждого сгенерированного объёма $b$ по каждой оси $p$ вырезается $S$ срезов. Индексы:
\[
  i_k = \left\lfloor \frac{(k + u_{b,p})\,L}{S} \right\rfloor,\qquad k=0,\dots,S-1,\quad u_{b,p}\sim\mathcal U[0,1),
\]
где сдвиг $u_{b,p}$ свой для каждого объёма и оси. Так срезы равномерно покрывают объём, а их положения
меняются от итерации к итерации. Одни и те же индексы берутся из пары $(x_t, x_{t+1})$, потому что
дискриминатор смотрит на согласованные пары.

%=====================================================================
\section{Диффузия (\code{diffusion.py})}

\subsection{Расписание шума}
Косинусное расписание (Nichol \& Dhariwal) с $T=8$ и $s=0{,}008$:
\[
  f(k) = \cos^2\!\left(\frac{k/T + s}{1+s}\cdot\frac{\pi}{2}\right),\qquad
  \abar_k = \operatorname{clip}\!\left(\frac{f(k)}{f(0)},\,10^{-5},\,1\right),\quad k=0,\dots,T,
\]
\[
  \alpha_k = \abar_k/\abar_{k-1},\qquad \beta_k = 1-\alpha_k,\qquad \alpha_0 = 1.
\]

\begin{center}
\small
\begin{tabular}{@{}lccccccccc@{}}
\toprule
$k$ & 0 & 1 & 2 & 3 & 4 & 5 & 6 & 7 & 8 \\
\midrule
$\abar_k$ & 1.000 & 0.958 & 0.847 & 0.684 & 0.494 & 0.304 & 0.144 & 0.038 & $10^{-5}$ \\
$\beta_k$ & 0 & 0.042 & 0.116 & 0.192 & 0.278 & 0.384 & 0.526 & 0.740 & 1.000 \\
\bottomrule
\end{tabular}
\end{center}

\subsection{Прямой процесс (реальные данные)}
Для реального кропа $x_0$ и шага $t\in\{0,\dots,T-1\}$ строится пара $(x_t, x_{t+1})$ (\code{q\_pair}):
\[
  x_t = \sqrt{\abar_t}\,x_0 + \sqrt{1-\abar_t}\,\varepsilon,\qquad
  x_{t+1} = \sqrt{\alpha_{t+1}}\,x_t + \sqrt{\beta_{t+1}}\,\varepsilon',\qquad
  \varepsilon,\varepsilon'\sim\mathcal N(0,I).
\]
При $t=0$ получается $x_0$ без шума ($\abar_0=1$).

\subsection{Обратный шаг (генерация)}
Генератор получает $x_{t+1}$, номер шага $t$ и латентный вектор $z$ и предсказывает чистый объём
$\hat x_0 = G_\theta(x_{t+1}, t, z)$. Следующее состояние сэмплируется из апостериорного распределения DDPM
(\code{posterior}):
\[
  q(x_t\mid x_{t+1}, \hat x_0) = \mathcal N\!\left(c_0(t)\,\hat x_0 + c_1(t)\,x_{t+1},\ \tilde\sigma_t^2 I\right),
\]
\[
  c_0(t) = \frac{\beta_{t+1}\sqrt{\abar_t}}{1-\abar_{t+1}},\qquad
  c_1(t) = \frac{(1-\abar_t)\sqrt{\alpha_{t+1}}}{1-\abar_{t+1}},\qquad
  \tilde\sigma_t^2 = \frac{\beta_{t+1}(1-\abar_t)}{1-\abar_{t+1}}.
\]
На последнем шаге ($t=0$) $\abar_0=1$, поэтому $c_0=1$, $c_1=0$, $\tilde\sigma_0=0$, и результат генерации
в точности равен $\hat x_0 = \tanh(\ell)$, где $\ell$ — выход сети до tanh.

\begin{center}
\small
\begin{tabular}{@{}lcccccccc@{}}
\toprule
$t$ & 0 & 1 & 2 & 3 & 4 & 5 & 6 & 7 \\
\midrule
$c_0(t)$ & 1.000 & 0.740 & 0.560 & 0.455 & 0.388 & 0.339 & 0.292 & 0.194 \\
$c_1(t)$ & 0 & 0.259 & 0.435 & 0.530 & 0.571 & 0.560 & 0.453 & 0.016 \\
$\tilde\sigma_t$ & 0 & 0.179 & 0.305 & 0.417 & 0.528 & 0.654 & 0.811 & 0.981 \\
\bottomrule
\end{tabular}
\end{center}

В статье обратный шаг записан упрощённо: $x_{t-1}=\sqrt{\abar_{t-1}}\hat x_0+\sqrt{1-\abar_{t-1}}\,\eta$.
В коде авторов используется стандартное апостериорное распределение, как выше. Мы следуем коду.

\subsection{Сэмплинг}
\begin{algorithm}[H]
\caption{Инференс (\code{Diffusion.sample})}
$x_T\sim\mathcal N(0,I)$ размера $n\times1\times L^3$;\quad $z\sim\mathcal N(0,I_{100})$ (один на всю траекторию)\;
\For{$t = T-1,\dots,0$}{
  $\hat x_0 \leftarrow G_\theta(x_{t+1}, t, z)$\;
  $x_t \leftarrow c_0(t)\hat x_0 + c_1(t)x_{t+1} + \mathbf{1}[t>0]\,\tilde\sigma_t\,\varepsilon$\;
}
\Return $x_0\in[-1,1]^{L^3}$ (при \code{logits=True} — логит $\ell$ последнего шага)\;
\end{algorithm}
При \code{logits=True} апостериорный шаг всё равно вызывается: он расходует одно обращение к генератору
случайных чисел, и следующие батчи получают тот же шум, что и при обычной генерации. Так результат с
порогом 0 совпадает бит в бит.

%=====================================================================
\section{Модели (\code{models.py})}

\subsection{3D-генератор $G_\theta(x_{t+1}, t, z)$}
\textbf{Условие.} Эмбеддинг шага и латентного вектора складываются:
\[
  e = E_t[t] + \mathrm{MLP}(z)\in\mathbb R^{128},\qquad z\sim\mathcal N(0,I_{100}),
\]
где $E_t$ — обучаемая таблица эмбеддингов ($T$ строк), а MLP состоит из трёх линейных слоёв с SiLU.

\textbf{Блок ResBlock3D с AdaGN.} Условие $e$ модулирует нормализованные признаки:
\[
  h = \mathrm{Conv}_{3^3}\big(\mathrm{SiLU}(\mathrm{GN}(x))\big),\qquad
  [\gamma,\ \delta] = W e + b,
\]
\[
  \mathrm{out} = \mathrm{Conv}_{3^3}\Big(\mathrm{SiLU}\big(\mathrm{GN}(h)\odot(1+\gamma) + \delta\big)\Big) + \mathrm{skip}(x),
\]
где $\mathrm{GN}$ — GroupNorm с 8 группами, а $\mathrm{skip}$ — свёртка $1^3$ при смене числа каналов.

\textbf{Архитектура (U-Net).}
\begin{center}
\small
\begin{tabular}{@{}llll@{}}
\toprule
Уровень & Разрешение & Каналы & Операции \\
\midrule
вход & $64^3$ & $1\to32$ & Conv $3^3$ \\
энкодер 1 & $64^3$ & 32 & ResBlock \\
энкодер 2 & $32^3$ & 64 & AvgPool, ResBlock \\
энкодер 3 & $16^3$ & 128 & AvgPool, ResBlock \\
энкодер 4 & $8^3$ & 128 & AvgPool, ResBlock \\
середина & $8^3$ & 128 & ResBlock \\
декодер 1 & $16^3$ & $128+128\to128$ & Upsample (nearest), concat, ResBlock \\
декодер 2 & $32^3$ & $128+64\to64$ & Upsample, concat, ResBlock \\
декодер 3 & $64^3$ & $64+32\to32$ & Upsample, concat, ResBlock \\
выход & $64^3$ & $32\to1$ & GN, SiLU, Conv $3^3$ $\to\ell$;\ \ $\hat x_0=\tanh(\ell)$ \\
\bottomrule
\end{tabular}
\end{center}
Всего 4{,}83\,M параметров. Сеть полностью свёрточная, поэтому на инференсе можно подавать объёмы любого
размера, кратного 8. Но обучена она на $64^3$, и на других размерах статистика смещается
(раздел~\ref{sec:limits}).

\subsection{2D-дискриминатор $D_\phi(x_t^{(p)}, x_{t+1}^{(p)}, t, p)$}
Один дискриминатор на все три плоскости. Вход — двухканальное изображение $64\times64$:
пара срезов $[x_t^{(p)}, x_{t+1}^{(p)}]$. Условие:
\[
  e = \big[E_t[t];\ E_p[p]\big]\in\mathbb R^{256}.
\]
Каждый слой получает условие как смещение по каналам:
\[
  h_l = \mathrm{LeakyReLU}_{0.2}\big(\mathrm{Conv}_l(h_{l-1}) + W_l e\big).
\]
Свёртки: $3\times3$ (64 канала, $64^2$), затем четыре свёртки $4\times4$ со страйдом 2: 128 ($32^2$),
256 ($16^2$), 512 ($8^2$), 512 ($4^2$). Голова — линейный слой $512\cdot16\to1$. Всего 7{,}34\,M параметров.

\textbf{Spectral normalization} стоит на всех свёртках и линейных слоях:
\[
  W_{\mathrm{SN}} = \frac{W}{\sigma_{\max}(W)}.
\]
$\sigma_{\max}$ оценивается степенным методом. Это ограничивает константу Липшица дискриминатора.

\textbf{Выход.} Логит обрезается и пропускается через tanh с температурой (раздел 3.3.1 статьи):
\[
  r = \operatorname{clip}(\mathrm{head}(h),\,-10,\,10),\qquad D = \tanh(r/\tau),\quad \tau=0{,}4.
\]
Поэтому выход дискриминатора лежит в $[-1,1]$, а минимум лосса D равен $2\,\mathrm{softplus}(-1)\approx0{,}63$.

%=====================================================================
\section{Обучение (\code{train.py})}

\subsection{Лоссы}
Используется non-saturating GAN loss. Реальные пары обозначены $x^{\mathrm r}=(x_t^{\mathrm r}, x_{t+1}^{\mathrm r})$,
сгенерированные — $x^{\mathrm g}=(x_t^{\mathrm g}, x_{t+1}^{\mathrm g})$:
\[
  \mathcal L_D = \E\big[\mathrm{softplus}(-D(x^{\mathrm r}))\big] + \E\big[\mathrm{softplus}(D(x^{\mathrm g}))\big],
  \qquad \mathrm{softplus}(u) = \log(1+e^u),
\]
\[
  \mathcal L_G = \E\big[\mathrm{softplus}(-D(x^{\mathrm g}))\big].
\]

\textbf{Lazy R1.} В каждой 16-й итерации (на всех 8 шагах $t$ этой итерации) к $\mathcal L_D$
добавляется штраф на градиент по реальным входам:
\[
  \mathcal L_{R1} = \frac{\gamma}{2}\,\E\Big[\big\|\nabla_{x_t^{\mathrm r}} D(x_t^{\mathrm r}, x_{t+1}^{\mathrm r}, t, p)\big\|_2^2\Big],\qquad \gamma = 10.
\]

\textbf{Прочие стабилизаторы.}
\begin{itemize}
  \item TTUR: Adam с $\beta=(0;\,0{,}9)$, $\eta_G = 2\cdot10^{-4}$, $\eta_D=10^{-5}$.
  \item Клиппинг нормы градиента генератора: $\|\nabla_\theta\|_2\le3$.
  \item Клиппинг логитов D в $[-10,10]$ и tanh с температурой (выше).
  \item bf16 autocast для прямых проходов. R1 считается в fp32.
\end{itemize}

\subsection{Low-pass curriculum}\label{sec:curr}
Эпоха — 100 итераций (в статье размер эпохи не указан).
\begin{center}
\small
\begin{tabular}{@{}lcc@{}}
\toprule
Итерации & $\sigma$ & Доля размытых кропов \\
\midrule
0--399 (эпохи 0--3) & 0{,}8 & 65\,\% \\
400--699 (эпохи 4--6) & 0{,}4 & 45\,\% \\
700 и далее & — & 0 \\
\bottomrule
\end{tabular}
\end{center}

\subsection{Итерация обучения}
Порядок повторяет код авторов (\code{train\_ddgan.py}): на каждом шаге цепочки сначала обновляется D,
затем G, после чего состояние отсоединяется от графа. В Algorithm 1 статьи сначала прогоняется вся
цепочка, затем делается одно обновление. Вариант авторского кода требует меньше памяти.

\begin{algorithm}[H]
\caption{Одна итерация (\code{Trainer.train\_iter})}
$x_T\sim\mathcal N(0,I)$, $B=2$ объёма $64^3$;\quad $z\sim\mathcal N(0,I_{100})$\;
\For{$t = T-1,\dots,0$}{
  $x_t \leftarrow \mathrm{posterior}\big(G_\theta(x_{t+1},t,z),\,x_{t+1},\,t\big)$ \tcp*{с графом}
  фейковые пары: $S=8$ стратифицированных срезов на ось из $(x_t, x_{t+1})$ $\Rightarrow$ $3\cdot S\cdot B=48$ пар\;
  реальные пары: 32 кропа, размытие по curriculum, $(x_t^{\mathrm r}, x_{t+1}^{\mathrm r})=\mathrm{q\_pair}(x_0^{\mathrm r}, t)$\;
  \textbf{шаг D:} $\mathcal L_D$ (фейки отсоединены) $+\ \mathcal L_{R1}$, если итерация кратна 16; Adam\;
  \textbf{шаг G:} $\mathcal L_G$ по тем же фейковым парам; клиппинг градиента; Adam\;
  $x_{t+1}\leftarrow \mathrm{stopgrad}(x_t)$\;
}
\end{algorithm}

\subsection{Гиперпараметры}
\begin{center}
\small
\begin{tabular}{@{}lll@{}}
\toprule
Параметр & Значение & Источник \\
\midrule
Размер объёма / кропа & $64^3$ / $64^2$ & статья, код авторов \\
Шагов диффузии $T$ & 8 & статья \\
Расписание & косинусное, $s=0{,}008$ & статья \\
Размерность $z$ & 100 & код авторов \\
Объёмов на итерацию $B$ & 2 & код авторов \\
Срезов на ось $S$ & 8 & наш выбор (у авторов 16) \\
Реальных кропов на шаг & 32 & наш выбор \\
$\eta_G$ / $\eta_D$ & $2\cdot10^{-4}$ / $10^{-5}$ & код авторов \\
Adam $\beta$ & $(0;\ 0{,}9)$ & код авторов \\
R1: $\gamma$, период & 10, каждые 16 итераций & статья \\
Клиппинг градиента G & 3 & статья \\
$\tau$ выхода D & 0{,}4 & статья \\
Итераций & 3000 (24\,000 обновлений D и G) & наш выбор \\
Seed обучения / генерации & 0 / 1 & — \\
\bottomrule
\end{tabular}
\end{center}

%=====================================================================
\section{Метрики (\code{metrics.py})}

\textbf{Доля фазы (пористость).} Для бинарного объёма $V$ из $N$ вокселей:
\[
  \phi = \frac1N\sum_{\mathbf x}\mathbf{1}[V(\mathbf x)=0].
\]

\textbf{Двухточечная функция.} Вероятность того, что обе точки на расстоянии $r$ вдоль оси $a$ лежат в
порах. Для индикатора $\chi(\mathbf x)=\mathbf{1}[V(\mathbf x)=0]$:
\[
  S_2^{(a)}(r) = \big\langle \chi(\mathbf x)\,\chi(\mathbf x + r\mathbf e_a)\big\rangle_{\mathbf x}.
\]
Считается через теорему Винера–Хинчина (периодическая граница) и усредняется по двум другим осям:
\[
  S_2^{(a)}(r) = \Big\langle \frac1{L}\,\mathcal F_a^{-1}\big[\,|\mathcal F_a[\chi]|^2\,\big](r)\Big\rangle,\qquad r=0,\dots,31.
\]
Свойства: $S_2(0)=\phi$, а при больших $r$ функция стремится к $\phi^2$.

\textbf{Эталон.} \code{random\_subvolumes} вырезает случайные кубы $64^3$ из керна для сравнения со
сгенерированными объёмами того же размера.

%=====================================================================
\section{Скрипты}

\subsection{\code{generate\_samples.py}}
Для каждого варианта: если нет \code{checkpoints/Variant\_XX.pt}, обучает модель (3000 итераций, seed 0).
Затем генерирует $N$ объёмов (seed 1, батчи по 4), бинаризует по порогу 0 и сохраняет
\code{samples.npy}, картинку со срезами и \code{stats.json} (время и пиковую память GPU).

\subsection{\code{calibrate\_threshold.py}}
Повторяет генерацию с тем же seed и проверяет, что порог 0 даёт в точности ранее сохранённый файл.
Затем подбирает порог по пористости \emph{2D-обучающих срезов}, то есть по информации, доступной методу:
\[
  \phi_{2D} = \frac13\sum_{p=0}^{2}\frac{1}{|I_p|}\sum_{\mathbf u}\mathbf{1}[I_p(\mathbf u)=-1],\qquad
  c = Q_{\phi_{2D}}\big(\{\ell_i(\mathbf x)\}_{i=1..N,\ \mathbf x}\big),
\]
где $Q_q$ — эмпирический $q$-квантиль логитов по всем $N$ объёмам. Порог один на вариант, поэтому
разброс между объёмами не подгоняется. Поры: $\ell<c$. Так как $\tanh$ монотонна, это равносильно порогу
$\tanh(c)$ на сером выходе. Но около 99{,}5\,\% вокселей насыщены до $\pm1$, и порог по серому значению
упирается в потолок, а по логиту — нет. Дополнительно скрипт считает пористость на гранях куба, чтобы
контролировать краевой артефакт.

\subsection{\code{view\_volume.py}}
Открывает \code{.npy} или \code{.mhd}. Режимы: 3D-изоповерхность выбранной фазы (marching cubes по уровню
0{,}5 после паддинга маски на 1 воксель, чтобы поверхность замыкалась на гранях, рендер в plotly),
листалка срезов (matplotlib) и экспорт в \code{.mhd}/\code{.raw.gz}.

%=====================================================================
\section{Результаты}

20 объёмов $64^3$ на вариант, RTX 5060 Ti 16 ГБ. $\phi_{2D}$ — пористость трёх обучающих срезов.

\begin{center}
\small
\begin{tabular}{@{}lcccccc@{}}
\toprule
Вариант & $\phi$ керна & $\phi_{2D}$ & $\phi$, порог 0 & $\phi$, калибровка & Обучение, мин & Генерация, с/объём \\
\midrule
V01 & 0.429 & 0.426 & $0.418\pm0.003$ & $0.426\pm0.003$ & $\approx47^*$ & 0.19 \\
V11 & 0.141 & 0.151 & $0.092\pm0.006$ & $0.151\pm0.012$ & 61.4 & 0.16 \\
V02 & 0.197 & 0.202 & $0.130\pm0.006$ & $0.202\pm0.008$ & 59.4 & 0.16 \\
V12 & 0.246 & 0.264 & $0.240\pm0.005$ & $0.264\pm0.005$ & $\approx81^*$ & 0.16 \\
V18 & 0.340 & 0.349 & $0.335\pm0.007$ & $0.349\pm0.008$ & 60.0 & 0.16 \\
\bottomrule
\end{tabular}
\end{center}
\noindent{\footnotesize $^*$ оценка по времени изменения файлов, замеры добавлены позже.
Пиковая память (PyTorch allocated): обучение $\approx2{,}5$ ГБ, генерация $\approx1{,}75$ ГБ.}

\subsection{Ограничения}\label{sec:limits}
\begin{itemize}
  \item \textbf{Заниженный разброс.} Стандартное отклонение пористости между объёмами в $\sim10$ раз меньше,
        чем у случайных кубов $64^3$ керна (V01: 0{,}003 против 0{,}034). Авторы отмечают то же.
  \item \textbf{Тесные породы.} На V11 и V02 при пороге 0 пористость занижена примерно на треть, проводимость
        сильно занижена (по измерениям коллеги). Калибровка порога исправляет пористость.
  \item \textbf{Краевой артефакт.} На гранях объёма пористость выше, чем внутри: после калибровки для V11
        крайний слой 0{,}23 против 0{,}15 внутри, для V02 0{,}30 против 0{,}20. У V12 то же уже при
        пороге 0. Генерация $80^3$ с вырезанием центра убирает артефакт, но сдвигает статистику
        (модель обучена на $64^3$).
  \item \textbf{Воспроизводимость.} Генерация из чекпойнта воспроизводится бит в бит. Обучение — нет:
        свёртки на GPU недетерминированы.
\end{itemize}

\subsection{Отличия от статьи и кода авторов}\label{sec:diff}
\begin{itemize}
  \item Генератор: компактный 3D U-Net с AdaGN (4{,}83\,M) вместо NCSN++ авторов.
  \item Дискриминатор написан по описанию в статье. tanh с температурой на выходе D взят из текста статьи,
        в коде авторов его нет.
  \item Порядок обновлений и апостериорное распределение — как в коде авторов, а не как в Algorithm 1 и
        формуле статьи.
  \item R1, клиппинг градиента и размытие — по статье (в коде авторов по умолчанию другие значения: R1
        выключен, клиппинг 12, размытие выключено).
  \item Размер эпохи (100 итераций), $S=8$, 32 реальных кропа и 3000 итераций — наш выбор.
  \item Обучающие данные: по одному центральному срезу на плоскость, как в статье (в конфигурации кода
        авторов по 5 срезов на ось).
\end{itemize}

%=====================================================================
\newpage
\appendix
\section{Листинг кода}

\subsection{\code{src/dimexdam/data.py}}
\codefile{../src/dimexdam/data.py}

\subsection{\code{src/dimexdam/diffusion.py}}
\codefile{../src/dimexdam/diffusion.py}

\subsection{\code{src/dimexdam/models.py}}
\codefile{../src/dimexdam/models.py}

\subsection{\code{src/dimexdam/train.py}}
\codefile{../src/dimexdam/train.py}

\subsection{\code{src/dimexdam/metrics.py}}
\codefile{../src/dimexdam/metrics.py}

\subsection{\code{scripts/generate\_samples.py}}
\codefile{../scripts/generate_samples.py}

\subsection{\code{scripts/calibrate\_threshold.py}}
\codefile{../scripts/calibrate_threshold.py}

\subsection{\code{scripts/view\_volume.py}}
\codefile{../scripts/view_volume.py}

\end{document}
